\chapter{On-Chain Trading}\label{ch:hf:on-chain-trading}

A large swap on a decentralised exchange moves its pool price away from the centralised exchanges. Within the same block, searchers compete to
arbitrage it back, and most of what they earn is bid away to the builder that includes their bundle. In this chapter's first-price auction, two
competing searchers leave the builder 93\% of an opportunity's value, ten leave it 98\%; the arbitrage itself, on a \$20 million pool over a
simulated week, nets its searchers \$4\,300 after gas. The liquidity provider on the other side of those trades, hedged on a centralised exchange,
loses \$26\,000 to them, the loss-versus-rebalancing a formula predicts to within 10\%, and keeps \$313\,000 of fees net of it.

\section{Exchange-to-chain arbitrage}

A constant-product pool (One Quant Book 3, chapter 20) quotes a price that moves only when someone trades with it. When the centralised
exchanges move, the pool is stale until an arbitrageur trades it back: the CEX--DEX arbitrage of One Quant Book 3, chapter 22. The optimal trade
has a closed form. For a pool of $x$ tokens and $y$ dollars with fee $1-g$, and a centralised price $p$ below the pool's price by more than the fee,
the arbitrageur sells tokens into the pool until the marginal output of the last one, $gxy/(x+g\,\Delta x)^2$, equals $p$:
\[
\Delta x^\ast=\frac{\sqrt{gxy/p}-x}{g},\qquad\text{profit}=\frac{g\,\Delta x^\ast y}{x+g\,\Delta x^\ast}-p\,\Delta x^\ast,
\]
and symmetrically with dollars when the pool is cheap (\cref{lst:hf:on-chain-trading:arb}). The trade leaves the pool within the fee of the
centralised price. The arbitrageur's hedge is a trade on the centralised exchange, which is not in the same block and not atomic: it carries price
risk between the two legs and the risk that the block does not include the on-chain leg.

\section{Atomic arbitrage and searching}

\begin{definition}[Atomic arbitrage]\label{def:hf:on-chain-trading:atomic}
\emph{Atomic arbitrage}\index{atomic arbitrage} is an arbitrage whose legs all execute in one on-chain transaction, typically across decentralised
exchanges and often with a flash loan, so that either every leg succeeds or the whole transaction reverts and only the fee for the failed
attempt is lost.
\end{definition}

Daian and co-authors documented the rise of arbitrage bots on decentralised exchanges, which, like high-frequency traders, paid high transaction
fees and optimised latency to front-run ordinary users' trades; they described priority gas auctions, in which bots bid up fees to obtain early
positions in the block, and showed that fees paid for ordering posed a risk to the consensus layer. That risk is what proposer--builder separation
(One Quant Book 3, chapter 22) reorganised: searchers now send bundles to builders, who assemble blocks, and proposers choose the most valuable
block.

\begin{omfigure}
\begin{tikzpicture}[font=\footnotesize,>=stealth,box/.style={draw,rounded corners,fill=omThm!8,minimum width=2.3cm,minimum height=0.75cm,align=center}]
\node[box] (u) at (0,1.2) {user swaps\\(mempool)};
\node[box] (c) at (0,-1.2) {centralised\\exchange price};
\node[box,fill=omDef!10] (s) at (3.6,0) {searchers\\($n$ bundles)};
\node[box] (b) at (8.0,0) {builder};
\node[box] (p) at (12.4,0) {proposer\\(block)};
\draw[->] (u) -- (s);
\draw[->] (c) -- (s);
\draw[->] (s) -- node[above=3pt] {bundle, bid} (b);
\draw[->] (b) -- node[below=3pt] {block, payment} (p);
\node[font=\scriptsize,align=center] at (5.8,-1.1) {first-price auction:\\the best bid is included};
\end{tikzpicture}

\omcaption{The path of an on-chain arbitrage: searchers see a user's swap or a centralised price move, build a bundle, and bid for its inclusion in
a first-price auction run by a builder, which assembles a block and pays the proposer for including it. Source: One Quant Book 3, chapter 22; this
chapter.}\label{fig:hf:on-chain-trading:pipeline}
\end{omfigure}

\section{Bundles, builders and bidding}

\begin{definition}[Bundle bid]\label{def:hf:on-chain-trading:bid}
A \emph{bundle bid}\index{bundle bid} is the payment a searcher attaches to a transaction bundle for its inclusion in a block, paid to the builder
(directly or through the transaction's priority fee) only if the bundle is included.
\end{definition}

When many searchers see the same opportunity, the bundle auction decides who gets it and at what price
(\cref{fig:hf:on-chain-trading:pipeline}). The chapter's model (\cref{lst:hf:on-chain-trading:auction}): an opportunity worth \$100 on public
information; each searcher's value is lower by its own cost (gas, inventory, the risk on the centralised leg), uniformly between 0\% and 20\%; bids
are sealed and the best is paid. With uniform private values the symmetric equilibrium bid shades the value toward the lowest by $1/n$. One
searcher pays nothing: it is a monopolist. With two, the builder receives 93\% of the winner's value; with five, 97\%; with twenty, 99\%
(\cref{fig:hf:on-chain-trading:auction}). What the searcher keeps is its cost advantage over the second-best, not the opportunity.

\begin{omfigure}
\begin{tikzpicture}
\begin{axis}[width=11cm,height=5.2cm,axis y line*=right,axis x line=none,xmode=log,log basis x=10,xmin=0.9,xmax=22,ymin=0,ymax=100,
  ylabel={\small winner keeps (\$ of \$100)},tick label style={font=\footnotesize},table/col sep=comma]
\addplot[omDef,thick,dashed,mark=square*,mark size=1.3pt] table[x=n,y=searcher]{figdata/hft/25-on-chain-trading/auction.csv};
\label{plot:hf:auckeep}
\end{axis}
\begin{axis}[width=11cm,height=5.2cm,axis y line*=left,xmode=log,log basis x=10,xmin=0.9,xmax=22,xtick={1,2,3,5,10,20},
  xticklabels={1,2,3,5,10,20},xlabel={\small competing searchers},ylabel={\small builder's share of the winner's value (\%)},ymin=0,
  ymax=105,tick label style={font=\footnotesize},grid=major,grid style={gray!25},table/col sep=comma,
  legend style={font=\scriptsize,at={(0.97,0.4)},anchor=east}]
\addplot[omThm,thick,mark=*,mark size=1.4pt] table[x=n,y=share]{figdata/hft/25-on-chain-trading/auction.csv};
\addlegendentry{builder's share (left)}
\addlegendimage{/pgfplots/refstyle=plot:hf:auckeep}\addlegendentry{winner keeps (right)}
\end{axis}
\end{tikzpicture}

\omcaption{A first-price bundle auction for a \$100 opportunity among searchers whose values are lower by their own costs, uniform between 0 and
20\%: the builder's share of the winning searcher's value, and what the winner keeps, against the number of competing searchers (log scale);
20\,000 auctions per point. Data: \texttt{hf\_onchain.auctions}.}\label{fig:hf:on-chain-trading:auction}
\end{omfigure}

\begin{dated}{2026-09}{The law and the block auction}\label{dat:hf:on-chain-trading:law}
In November 2025 a federal judge in Manhattan declared a mistrial in the case of Anton and James Peraire-Bueno, charged with wire fraud,
conspiracy to commit wire fraud and conspiracy to commit money laundering over an alleged exploit of the MEV-Boost software that netted \$25 million
in twelve seconds; the jury, reports said, could not agree on how to apply the law. Prosecutors could retry the case.
\end{dated}

\section{Liquidations on lending protocols}

Lending protocols (One Quant Book 3, chapter 23) let anyone repay an under-collateralised borrower's debt and take its collateral at a bonus.
Liquidating \$1 million of debt at a 5\% bonus is worth \$50\,000 less gas; many searchers watch every position's health factor, and the
opportunity goes to auction exactly like an arbitrage. Flash loans let a searcher liquidate without capital, borrowing, repaying the debt, selling
the collateral and repaying the loan in one atomic transaction.

\section{Just-in-time and hedged liquidity provision}

A liquidity provider in the pool takes the other side of every arbitrage. Milionis, Moallemi, Roughgarden and Zhang decomposed its return into a
market exposure and a microstructure part: fees earned less losses to arbitrageurs. For a constant-product pool worth $V$, with the token's
volatility $\sigma$, the loss-versus-rebalancing runs at $\sigma^2V/8$ per unit time; a provider that hedges the pool's token exposure on a
centralised exchange keeps fees less that rate.

\begin{definition}[Hedged liquidity provision]\label{def:hf:on-chain-trading:hedged}
\emph{Hedged liquidity provision}\index{hedged liquidity provision} is supplying liquidity to an automated market maker while hedging the
position's exposure to the token's price on another venue, so that the provider's result is the pool's fees less its loss to arbitrageurs, with
the token's price risk removed.
\end{definition}

The chapter's pool (\cref{lst:hf:on-chain-trading:pool}): \$20 million at a 0.3\% fee, a token at \$3\,000 with 4\% daily volatility, 12-second blocks
for a week, \$2\,000 of noise swaps a block in random directions, arbitrage after each whenever it beats \$5 of gas. Over three simulated weeks the
loss-versus-rebalancing is \$25\,800 to \$29\,500 against the formula's \$28\,000; the arbitrageurs also take back the price impact the noise traders
paid, \$20\,000 to \$22\,500 a week, which is theirs by construction. The hedged provider earns \$339\,000 of fees and nets \$313\,000.

\begin{omfigure}
\begin{tikzpicture}
\begin{axis}[width=11cm,height=5.2cm,xlabel={\small pool fee (basis points)},ylabel={\small \$ thousand over the week},xmode=log,log basis x=10,
  xtick={5,10,30,100},xticklabels={5,10,30,100},xmin=4,xmax=120,ymode=log,log basis y=10,ymin=10,ymax=1500,
  tick label style={font=\footnotesize},grid=major,grid style={gray!25},table/col sep=comma,
  legend style={font=\scriptsize,at={(0.03,0.97)},anchor=north west}]
\addplot[omThm,thick,mark=*,mark size=1.4pt] table[x=fee,y=fees]{figdata/hft/25-on-chain-trading/lp.csv};
\addplot[omDef,thick,dashed,mark=square*,mark size=1.3pt] table[x=fee,y=lvr]{figdata/hft/25-on-chain-trading/lp.csv};
\addplot[gray,thick,dotted,mark=triangle*,mark size=1.3pt] table[x=fee,y=net]{figdata/hft/25-on-chain-trading/lp.csv};
\legend{fees earned,loss-versus-rebalancing,hedged net}
\end{axis}
\end{tikzpicture}

\omcaption{The hedged liquidity provider's week on a \$20 million pool against its fee tier (log scales): fees earned, loss-versus-rebalancing and
the net; the noise traders' volume is held fixed, which flatters high fees (real noise volume falls with the fee). Data:
\texttt{hf\_onchain.by\_fee}.}\label{fig:hf:on-chain-trading:lp}
\end{omfigure}

Loss-versus-rebalancing barely depends on the fee (\cref{fig:hf:on-chain-trading:lp}): it is set by volatility and the pool's size. The fee decides
the income, through the noise volume the pool attracts; the model holds that volume fixed, which is exactly the assumption a provider must not
make when it chooses a fee tier. Just-in-time liquidity (One Quant Book 3, chapter 20) is the extreme of hedged provision: a searcher adds
concentrated liquidity around a large swap it has seen in the mempool, earns the swap's fee, and removes the liquidity in the same block.

\section{Sandwiching and the law}

A sandwich attack (One Quant Book 3, chapter 22) buys ahead of a victim's swap and sells after it, taking the price impact the victim was willing to
tolerate. On the chapter's pool, a \$100\,000 purchase with a 1\% slippage tolerance can be preceded by a \$50\,900 front-run that leaves the victim
0.33 tokens, about \$990, worse off, exactly its tolerance; the attacker's gross gain is \$700, of which, in a competitive block auction, 90\% goes
to the builder and \$10 to gas, leaving \$60 (Book 3's \texttt{firm.sandwich}). The victim's loss is the attack's revenue; the builder's share is
paid out of it. Whether that is fraud, market manipulation or neither has not been settled in court: the one prosecution that reached a jury, over
an exploit that targeted sandwich bots themselves, ended in a mistrial (dated box). This chapter describes sandwiching only to price what its
victims lose; it is not a strategy file.

\section{Strategy files}

\begin{strategyfile}{CEX-DEX arbitrage}\label{strat:hf:on-chain-trading:cexdex}
\sfield{Who pays you, and why} Liquidity providers in stale pools, through loss-versus-rebalancing.
\sfield{Instruments and venues} Constant-product and concentrated pools; the centralised exchanges' books.
\sfield{Signal} The pool's price against the centralised price, beyond the fee and gas.
\sfield{Sizing and execution} The closed-form size; the on-chain leg in a bundle, the centralised leg as a hedge.
\sfield{Costs} Gas, the \omterm{def:hf:on-chain-trading:bid}{bundle bid} (93\% to 99\% of the value against competitors), the centralised leg's risk and fees.
\sfield{How it dies} Competition in the bundle auction; in the model the week's arbitrage nets its searchers \$4\,300 after gas, before bids.
\sfield{Horizon, capacity, infrastructure} One block; inventory on both sides; builder connections.
\sfield{Backtest honestly} Block-by-block pool states, the centralised price at each block's time, bids that would have won.
\sfield{Sources} Milionis, Moallemi, Roughgarden and Zhang; One Quant Book 3, chapter 22.
\end{strategyfile}

\begin{strategyfile}{Atomic DEX-DEX arbitrage}\label{strat:hf:on-chain-trading:atomic}
\sfield{Who pays you, and why} Traders whose swaps leave two pools at different prices.
\sfield{Instruments and venues} Pools of the same pair, or cycles of pairs, on decentralised exchanges.
\sfield{Signal} A cycle whose output exceeds its input after fees.
\sfield{Sizing and execution} One atomic transaction, often flash-funded; reverts if the price moved.
\sfield{Costs} Gas, including for reverted attempts; the \omterm{def:hf:on-chain-trading:bid}{bundle bid}.
\sfield{How it dies} The priority gas auctions Daian and co-authors documented, and their successor, the builder auction.
\sfield{Horizon, capacity, infrastructure} One block; mempool and state simulation.
\sfield{Backtest honestly} Simulate against the chain's state at the block, including competing bundles.
\sfield{Sources} Daian et al.\ (2019, Flash Boys 2.0).
\end{strategyfile}

\begin{strategyfile}{Lending-protocol liquidation searching}\label{strat:hf:on-chain-trading:liquidation}
\sfield{Who pays you, and why} Borrowers whose collateral falls, through the protocol's liquidation bonus.
\sfield{Instruments and venues} Lending protocols; the collateral's markets to sell it.
\sfield{Signal} Health factors below one, from on-chain positions and prices.
\sfield{Sizing and execution} Repay the debt, take the collateral at the bonus, sell it; flash loans if needed; bid for inclusion.
\sfield{Costs} Gas, the bid, the collateral's sale in a falling market.
\sfield{How it dies} Crowded auctions in crashes, when gas and the collateral's slippage peak together.
\sfield{Horizon, capacity, infrastructure} One block; oracle and position monitoring.
\sfield{Backtest honestly} Oracle updates as they were, competing liquidators, sale prices in the crash.
\sfield{Sources} One Quant Book 3, chapter 23.
\end{strategyfile}

\begin{strategyfile}{Just-in-time liquidity provision}\label{strat:hf:on-chain-trading:jit}
\sfield{Who pays you, and why} Large swappers, through the fee on their swap.
\sfield{Instruments and venues} Concentrated-liquidity pools.
\sfield{Signal} A large pending swap in the mempool or an order-flow auction.
\sfield{Sizing and execution} Add liquidity around the price for the swap, remove it after, in one bundle; hedge the inventory the swap leaves.
\sfield{Costs} Gas, the \omterm{def:hf:on-chain-trading:bid}{bundle bid}, the hedge.
\sfield{How it dies} Private order flow that hides swaps until they are included.
\sfield{Horizon, capacity, infrastructure} One block.
\sfield{Backtest honestly} Only swaps that were visible before inclusion.
\sfield{Sources} One Quant Book 3, chapter 20.
\end{strategyfile}

\begin{strategyfile}{Bundle bidding in builder auctions}\label{strat:hf:on-chain-trading:bidding}
\sfield{Who pays you, and why} Nobody pays the bidder: the bid decides how much of an opportunity the searcher keeps.
\sfield{Instruments and venues} Builder auctions for bundles.
\sfield{Signal} The opportunity's value, the searcher's own cost, the number of competitors.
\sfield{Sizing and execution} Shade the bid by the equilibrium: toward the lowest value by $1/n$ with uniform costs.
\sfield{Costs} Losing auctions; overbidding.
\sfield{How it dies} More competitors: the winner keeps \$6.7 of \$100 against one rival and \$1.0 against nineteen; and exploits of the auction's
software, the subject of the Peraire-Bueno prosecution (dated box).
\sfield{Horizon, capacity, infrastructure} One block.
\sfield{Backtest honestly} Competitors' bids from the chain, not the searcher's own.
\sfield{Sources} This chapter; the dated box.
\end{strategyfile}

\begin{strategyfile}{Hedged concentrated-liquidity market making}\label{strat:hf:on-chain-trading:hedgedlp}
\sfield{Who pays you, and why} Swappers, through fees; less what arbitrageurs take.
\sfield{Instruments and venues} Concentrated-liquidity pools; a centralised exchange for the hedge.
\sfield{Signal} Fee income against the loss-versus-rebalancing rate $\sigma^2V/8$.
\sfield{Sizing and execution} Provide where fees beat the rate; hedge the pool's delta each block or beyond a band.
\sfield{Costs} Loss-versus-rebalancing (\$26\,000 a week on \$20 million at 4\% daily volatility), hedging, gas.
\sfield{How it dies} Volatility spikes, which raise the loss with the square of volatility while fees follow volume.
\sfield{Horizon, capacity, infrastructure} Days to months.
\sfield{Backtest honestly} Noise volume that responds to the fee tier.
\sfield{Sources} Milionis, Moallemi, Roughgarden and Zhang.
\end{strategyfile}

\section{Tutorial: bid away to the builder}\label{tut:hf:on-chain-trading:tut}

\begin{tutorial}
\textbf{Goal.} Price the arbitrage, the bundle auction and the hedged liquidity provider's week on a simulated chain. \textbf{End state:} the
figures and the numbers in the text.
\begin{tutsteps}
\item \textbf{The optimal arbitrage} in closed form, checked against \texttt{firm.amm}'s integer swap.
\omcode{code/firm/searcher/firm_searcher.py}{38}{51}{Sell into a dear pool or buy from a cheap one until the marginal unit earns the centralised price.}\label{lst:hf:on-chain-trading:arb}
\item \textbf{The auction} among $n$ searchers.
\omcode{code/firm/searcher/firm_searcher.py}{66}{75}{Uniform private values and the first-price equilibrium bid; the builder's share.}\label{lst:hf:on-chain-trading:auction}
\item \textbf{The pool's week}: noise swaps, arbitrage after each block, the hedged provider's account.
\omcode{code/firm/searcher/firm_searcher.py}{103}{131}{Each block's noise swap and arbitrage, valued at the centralised price.}\label{lst:hf:on-chain-trading:pool}
\end{tutsteps}
\textbf{What to change next.} Let noise volume fall with the fee; add a second pool and search atomic cycles; put the searchers' latency to the
centralised exchange into their costs.
\end{tutorial}

\section{Build: the searcher}\label{bld:hf:on-chain-trading:searcher}

\begin{build}
\bfield{Purpose} Price on-chain arbitrage and liquidations, the bundle auction that allocates them, and a hedged provider's loss to them.
\bfield{Interface} \texttt{optimal\_arb(x, y, p, fee)}, \texttt{arb\_exact}, \texttt{auction(value, n, spread, seed, trials)}, \texttt{Chain(blocks,
seed, sigma\_day, p0, noise\_usd)}, \texttt{run\_pool(chain, value0, fee, gas\_usd)}, \texttt{lvr\_rate}, \texttt{liquidation\_value},
\texttt{gas\_cost}. Built on Book 3's \texttt{firm.amm} and \texttt{firm.gasfee}; the sandwich arithmetic is Book 3's \texttt{firm.sandwich}.
\bfield{Rules} Dollars; fees as fractions; one arbitrage a block after the noise swap.
\bfield{Acceptance tests} \texttt{code/firm/searcher/tests/}: the arbitrage's first-order condition and its optimality against neighbouring sizes; no
arbitrage inside the fee band; agreement with the integer swap; the auction's split by hand; loss-versus-rebalancing within 40\% of the formula over
three days; the hedged account adds up.
\bfield{Stretch} Atomic cycles; order-flow auctions; latency-dependent searcher costs.
\end{build}

\begin{omsources}
\item P.~Daian et al., Flash Boys 2.0: frontrunning, transaction reordering, and consensus instability in decentralized exchanges,
arXiv:1904.05234, 2019 (IEEE Symposium on Security and Privacy, 2020).
\item J.~Milionis, C.~C.~Moallemi, T.~Roughgarden, A.~L.~Zhang, Automated market making and loss-versus-rebalancing, arXiv:2208.06046.
\item The Block, report of the mistrial in \emph{United States v.\ Peraire-Bueno}, 8 November 2025.
\end{omsources}

\section{Exercises}

\begin{exercise}[$\star$]\label{exo:hf:on-chain-trading:1}
A pool holds 1\,000 tokens and \$3\,100\,000; the centralised price is \$3\,000 and the fee 0.3\%. How many tokens should the arbitrageur sell into
the pool, and for what profit?
\end{exercise}

\begin{exercise}[$\star$]\label{exo:hf:on-chain-trading:2}
What is the loss-versus-rebalancing rate of a \$20 million constant-product pool at 4\% daily volatility, per day and per week?
\end{exercise}

\begin{exercise}[$\star$]\label{exo:hf:on-chain-trading:3}
Liquidating \$1 million of debt at a 5\% bonus costs \$20 of gas. What is it worth?
\end{exercise}

\begin{exercise}[$\star\star$]\label{exo:hf:on-chain-trading:4}
Why does the builder's share rise toward 100\% as searchers are added?
\end{exercise}

\begin{exercise}[$\star\star$]\label{exo:hf:on-chain-trading:5}
Why does loss-versus-rebalancing barely depend on the pool's fee in the model?
\end{exercise}

\begin{exercise}[$\star\star$]\label{exo:hf:on-chain-trading:6}
Why is a CEX--DEX arbitrage not atomic, and what risk does that leave?
\end{exercise}

\begin{exercise}[$\star\star\star$]\label{exo:hf:on-chain-trading:7}
\emph{Coding.} Double the token's volatility in \texttt{firm.searcher.Chain}. What happens to loss-versus-rebalancing and to the hedged provider's net?
\end{exercise}

\begin{exercise}[$\star\star\star$]\label{exo:hf:on-chain-trading:8}
\emph{Find the flaw.} ``Arbitrage opportunities on our pools were worth \$11\,400 last week; our searcher will make that.''
\end{exercise}

\section{Problem: Bid Away to the Builder}

\begin{problem}[{Weekend problem --- bid away to the builder}]\label{pb:hf:on-chain-trading:1}
A firm runs searchers and a hedged liquidity position on a decentralised exchange.

\medskip\noindent\textbf{Part I --- Arbitrage.}
\begin{enumerate}
\item Derive the optimal CEX--DEX trade size.
\item Define \omterm{def:hf:on-chain-trading:atomic}{atomic arbitrage} and contrast it with CEX--DEX arbitrage.
\item What did Daian and co-authors document?
\item Describe the path of a bundle from searcher to block.
\end{enumerate}

\medskip\noindent\textbf{Part II --- Auctions.}
\begin{enumerate}[resume]
\item Define a \omterm{def:hf:on-chain-trading:bid}{bundle bid}.
\item Give the equilibrium bid with uniform private values, and the builder's share for 2, 5 and 20 searchers.
\item What does the winner keep, and why?
\item How is a liquidation auctioned?
\end{enumerate}

\medskip\noindent\textbf{Part III --- Liquidity.}
\begin{enumerate}[resume]
\item Define \omterm{def:hf:on-chain-trading:hedged}{hedged liquidity provision} and the loss-versus-rebalancing rate.
\item Compare the simulated loss with the formula.
\item What else do arbitrageurs take from the pool, and who paid for it?
\item Why is the model's fee comparison flattering to high fees?
\end{enumerate}

\medskip\noindent\textbf{Part IV --- The verdict.}
\begin{enumerate}[resume]
\item State the \emph{named result}: the share of arbitrage profit paid to builders as the number of competing searchers grows, and the hedged
liquidity provider's net return after loss-versus-rebalancing.
\item What does a searcher need to keep more than 1\%?
\item What does the sandwich example show about who pays whom?
\item Summarise the dated box.
\item Which strategy file is most exposed to private order flow?
\item How would you choose a pool's fee tier?
\item What would an order-flow auction change for the swapper?
\item In one sentence: who earns an on-chain arbitrage?
\end{enumerate}
\end{problem}

\section{Interview questions}

\begin{interviewq}[$\star$ \iqroles{trader}]\label{iq:hf:on-chain-trading:1}
Ether moves 1\% on the centralised exchanges. What happens to a constant-product pool of ether in the next block, and who profits?
\end{interviewq}

\begin{interviewq}[$\star\star$ \iqroles{researcher}]\label{iq:hf:on-chain-trading:2}
Derive the loss-versus-rebalancing rate of a constant-product pool.
\end{interviewq}

\begin{interviewq}[$\star\star$ \iqroles{developer}]\label{iq:hf:on-chain-trading:3}
Design a searcher that finds, simulates and bids for arbitrage bundles within one block time.
\end{interviewq}

\begin{interviewq}[$\star\star$ \iqroles{risk}]\label{iq:hf:on-chain-trading:4}
Your CEX--DEX arbitrage's on-chain leg was not included but the centralised hedge was filled. What is your position, and what should the system do?
\end{interviewq}

\begin{interviewq}[$\star\star$ \iqroles{trader}]\label{iq:hf:on-chain-trading:5}
How would you bid against a competitor whose costs are lower than yours?
\end{interviewq}

\begin{interviewq}[$\star\star\star$ \iqroles{researcher}]\label{iq:hf:on-chain-trading:6}
Show that with $n$ bidders whose values are uniform on $[a,b]$, bidding $a+\frac{n-1}{n}(v-a)$ is a symmetric equilibrium of the first-price
auction.
\end{interviewq}
